\documentclass[13pt, a4paper, oneside]{scrreprt}

\usepackage{scrhack}
\usepackage[utf8]{inputenc}
\usepackage[T1]{fontenc}
\usepackage[polish]{babel}
\usepackage{csquotes}
\usepackage{mathpazo}
\usepackage[scaled=0.95]{helvet}
\usepackage{courier}
\usepackage{microtype}
\usepackage{setspace}
\usepackage[margin=2.5cm]{geometry}
\usepackage[automark,headsepline]{scrlayer-scrpage}
\usepackage{xcolor}
\usepackage{graphicx}
\usepackage[newfloat]{minted}
\usepackage{caption}
\usepackage{subcaption}
\usepackage[style=numeric, sorting=nty]{biblatex}
\usepackage{booktabs}
\usepackage{float}
\usepackage{amsmath}
\usepackage{siunitx}
\sisetup{output-decimal-marker={,}}
\usepackage{array}
\usepackage{tabularx}
\usepackage{placeins}

\definecolor{myblue}{RGB}{20, 80, 150}

\usepackage[
    colorlinks=true,
    linkcolor=myblue,
    citecolor=myblue,
    urlcolor=myblue,
    bookmarks=true
]{hyperref}

\onehalfspacing
\pagestyle{scrheadings}
\setlength{\parskip}{1em}

\SetupFloatingEnvironment{listing}{
    fileext=lol,
    listname={Spis kodów},
    name=Listing,
    placement=htbp
}

\setcounter{biburlnumpenalty}{100}
\setcounter{biburlucpenalty}{100}
\setcounter{biburllcpenalty}{100}

\addbibresource{bibliografia.bib}

\title{Predykcja zainteresowania postami w social media z użyciem metod NLP}
\author{Alicja Borek \and Dobrawa Rumszewicz \and Patryk Kowalczyk}
\date{Czerwiec 2026}

% Wykresy eksportowane z notebooków należy umieścić w katalogu docs/figures/.
% Makro poniżej pokaże wykres, jeżeli plik istnieje. W przeciwnym razie
% zostawi czytelny placeholder z nazwą oczekiwanego pliku.
\newcommand{\autoplot}[1]{%
    \IfFileExists{figures/#1}{%
        \includegraphics[width=\linewidth]{figures/#1}%
    }{%
        \fbox{%
            \parbox[c][0.23\textheight][c]{0.92\linewidth}{%
                \centering\small
                Miejsce na wykres wygenerowany w notebooku.\\[0.5em]
                Wstaw plik graficzny do katalogu \texttt{figures}.%
            }%
        }%
    }%
}

\newcommand{\resultplaceholder}[1]{%
    \textcolor{myblue}{\textbf{[Uzupełnić po ponownym uruchomieniu notebooków: #1]}}%
}

\begin{document}

\maketitle

\tableofcontents
\newpage

\chapter{Wprowadzenie}

Jedną z podstawowych form reakcji użytkowników serwisu Twitter, obecnie działającego pod nazwą X,
było przekazanie dalej cudzego wpisu. W analizowanym zbiorze jest to zapisane jako liczba
retweetów. W tym projekcie liczba retweetów została przyjęta jako miara zainteresowania wpisem.
Celem jest sprawdzenie, czy z formy i treści wpisu można przewidzieć, czy znalazł się on w
grupie tweetów o niskim, średnim albo wysokim zainteresowaniu.

Temat wymaga zbudowania modelu predykcyjnego dla tweetów Donalda Trumpa,
wykorzystania reprezentacji słów występujących w tekście oraz przeprowadzenia analizy SHAP.
NLP \emph{Natural Language Processing} oznacza przetwarzanie języka naturalnego,
czyli przekształcanie tekstu napisanego przez człowieka
w postać liczbową, którą może wykorzystywać algorytm uczenia maszynowego.

Głównym pytaniem jest: czy słowa, frazy oraz proste cechy stylistyczne tweetów pozwalają
trafnie przewidywać poziom liczby retweetów? Drugie pytanie dotyczy interpretacji modelu:
które elementy wpisu najsilniej wpływają na decyzję modelu dotyczącą
klasy \emph{Low}, \emph{Medium} albo \emph{High}? Odpowiedź na drugie pytanie uzyskano
przy użyciu metody SHAP, czyli \emph{SHapley Additive exPlanations}.

W literaturze dotyczącej konta \texttt{@realDonaldTrump} zwraca się uwagę zarówno na
tematykę, jak i na styl wypowiedzi. Clarke i Grieve opisali zmienność stylu tweetów
publikowanych na tym koncie i wyróżnili style konwersacyjny, kampanijny, zaangażowany
oraz doradczy \cite{clarke2019}. Savoy wskazał natomiast na rolę charakterystycznego
słownictwa, negatywnego tonu, ocen kategorycznych oraz wzmacniających przymiotników
\cite{savoy2022}. Te obserwacje uzasadniają wykorzystanie nie tylko pojedynczych słów,
ale również cech takich jak liczba wykrzykników, wzmianki o innych użytkownikach,
długość tekstu, liczby, czy wskaźniki czytelności. Materiały użyte użyte jako inspiracja
do wyboru analizowanych cech, zwracają uwagę na podobne zjawiska stylistyczne,
między innymi powtarzalność słów, krótkie zdania, superlatywy oraz wzmianki i hashtagi
\cite{textInspector2022,theConversation2018,healy2018}.

\chapter{Dane i wstępna analiza eksploracyjna}

\section{Źródło i wybór zbioru danych}

W projekcie wykorzystano zbiór \emph{Trump Tweets} udostępniony
na platformie Kaggle \cite{kaggleTrumpTweets}. Po rozpakowaniu danych dostępne
były dwa pliki: \texttt{realdonaldtrump.csv} oraz \texttt{trumptweets.csv}.
Pierwszy plik zawierał \num{43352} rekordy, natomiast drugi \num{41122}.
Porównanie identyfikatorów tweetów wykazało, że oba pliki mają \num{41104} wspólnych
wpisów. W pliku \texttt{realdonaldtrump.csv} znajdowało się dodatkowo \num{2248} tweetów,
a w drugim pliku tylko \num{18} tweetów nieobecnych w pierwszym. Z tego powodu do dalszej
analizy wybrano plik \texttt{realdonaldtrump.csv}.

Każdy rekord zawierał identyfikator wpisu, odnośnik, pełną treść, datę publikacji,
liczbę retweetów, liczbę polubień, a w części przypadków również informacje o wzmiankach
i hashtagach. Kolumna \texttt{content} była podstawowym źródłem tekstu.
Kolumna \texttt{retweets} została wykorzystana do zbudowania zmiennej docelowej.
Kolumna \texttt{favorites} została wykorzystana wyłącznie w analizie eksploracyjnej i nie znalazła
się w końcowej macierzy cech. Włączenie liczby polubień do modelu przewidującego retweety
byłoby dużym pominięciem analizy tekstu, ponieważ jak wynika z macierzy korelacji \ref{fig:feature-correlation}
kolumny \texttt{retweets} i \texttt{favorites} są mocno powiązane. Model mógłby wtedy uzyskiwać
wysoki wynik dzięki informacjom dostępnym z pierwszych polubień, a nie dzięki analizie języka tweeta.

\section{Rozkład retweetów i definicja klas}

Liczba retweetów ma bardzo nierówny rozkład. Mediana wynosiła 396{,}5,
średnia wynosiła około \num{6265}, a maksymalna wartość osiągała \num{302269}.
Duża różnica między średnią i medianą wskazuje, że niewielka liczba bardzo
popularnych tweetów silnie podnosi średnią, co potwierdza wykres \ref{fig:eda-retweets-length}.
Taki rozkład utrudniałby bezpośrednią regresję liczby retweetów, dlatego stworzono
klasyfikację trzech poziomów zainteresowania.

Granice klas wyznaczono na podstawie kwantyli rozkładu retweetów. Tweet
z nie więcej niż 53 retweetami otrzymał etykietę \emph{Low}. Tweet z liczbą
retweetów od 54 do 4558 otrzymał etykietę \emph{Medium}. Tweet z liczbą retweetów
większą niż 4558 otrzymał etykietę \emph{High}, co przedstawiono w tabeli \ref{fig:eda-classes}.
Zastosowanie dwóch kwantyli, odpowiadających około jednej trzeciej i dwóm trzecim obserwacji,
prowadzi do klas o bardzo zbliżonej liczności. Dzięki temu model nie powinien uzyskiwać dobrego
wyniku wyłącznie przez wybieranie najczęstszej klasy.

\begin{table}[H]
\centering
\caption{Podstawowe informacje o danych i zadaniu klasyfikacyjnym.}
\label{tab:dataset-summary}
\begin{tabularx}{\textwidth}{>{\raggedright\arraybackslash}p{0.46\textwidth}X}
\toprule
Element & Wartość \\
\midrule
Wybrany plik źródłowy & \texttt{realdonaldtrump.csv} \\
Liczba tweetów & \num{43352} \\
Liczba klas & 3: \emph{High}, \emph{Low}, \emph{Medium} \\
Granica klasy \emph{Low} & Maksymalnie 53 retweety \\
Granica klasy \emph{Medium} & Od 54 do 4558 retweetów \\
Granica klasy \emph{High} & Powyżej 4558 retweetów \\
Liczba przykładów treningowych & \num{34681} \\
Liczba przykładów testowych & \num{8671} \\
\bottomrule
\end{tabularx}
\end{table}

\begin{figure}[H]
\centering
\begin{subfigure}[t]{0.48\textwidth}
\centering
\autoplot{eda_retweets_and_length.png}
\caption{Rozkład liczby retweetów oraz długości tweetów.}
\label{fig:eda-retweets-length}
\end{subfigure}
\hfill
\begin{subfigure}[t]{0.48\textwidth}
\centering
\autoplot{eda_engagement_classes.png}
\caption{Granice i liczności klas zainteresowania.}
\label{fig:eda-classes}
\end{subfigure}
\caption{Wykresy z pierwszego etapu analizy eksploracyjnej.}
\label{fig:eda-main}
\end{figure}

\section{Zmiana zainteresowania w czasie}

W analizie eksploracyjnej data publikacji została przekształcona do formatu daty,
a następnie dla każdego roku obliczono medianę liczby retweetów oraz liczbę
opublikowanych wpisów. Najwyższą medianę retweetów zaobserwowano w roku 2020, kiedy
wynosiła ona \num{19279}. Z danych widać, że popularność konta konsekwentnie rosła
z czasem, co widoczne jest na wykresie \ref{fig:eda-time}.
Zjawisko to jest istotne interpretacyjnie, ponieważ wysoka
liczba retweetów zależy nie tylko od słów użytych w tweecie, lecz również od okresu
jego publikacji, bieżących wydarzeń i liczby obserwujących konto.

Zmienna czasu nie została jednak dodana do głównej macierzy cech.
Celem projektu było sprawdzenie ile informacji można wyciągnąć z tekstu i jego cech językowych.
Dodanie roku, miesiąca albo godziny mogłoby poprawić wynik predykcyjny, ale
osłabiłoby odpowiedź na pytanie, czy sam styl i treść wpisu wyjaśniają zainteresowanie.
Z tego powodu analiza czasu została potraktowana jako element rozpoznania danych,
a nie jako główne źródło predykcji.

\begin{figure}[H]
\centering
\begin{subfigure}[t]{0.58\textwidth}
\centering
\autoplot{eda_temporal_analysis.png}
\caption{Mediana retweetów i liczba tweetów według roku.}
\label{fig:eda-time}
\end{subfigure}
\hfill
\begin{subfigure}[t]{0.38\textwidth}
\centering
\autoplot{eda_early_correlation.png}
\caption{Wczesna korelacja Pearsona między retweetami, długością wpisu i liczbą polubień.}
\label{fig:eda-correlation}
\end{subfigure}
\caption{Analiza zmian zainteresowania w czasie i wstępnych zależności liczbowych.}
\label{fig:eda-time-correlation}
\end{figure}

\chapter{Przygotowanie tekstu}

\section{Czyszczenie tweetów}

Tweety zawierają elementy, które mogą utrudniać analizę słów.
Należą do nich adresy internetowe, odnośniki do obrazów, emotikony, wielokrotne odstępy
i znaki specjalne. W pierwszym etapie każdy tweet został zamieniony na małe litery,
a następnie usunięto adresy rozpoczynające się od \texttt{http},
\texttt{www} oraz odnośniki \texttt{pic.twitter.com}. Zachowano cyfry,
ponieważ mogą one wskazywać na lata, liczby głosów, wielkości ekonomiczne
lub inne fakty, które mają znaczenie dla zainteresowania.

Jeśli chodzi o wzmianki (ang. \textit{mentions}) i hashtagi to postanowiono zachować o nich informację
więc zamieniono je do takiego formatu aby znaki \texttt{@} i \texttt{\#} nie zostały potraktowane jak znaki specjalne i usunięte.
Wzmianka o koncie jest przekształcana do tokenu w postaci
\texttt{mentions\_nazwa\_konta}, na przykład \texttt{@CNN} może zostać
zapisana jako \texttt{mentions\_cnn}. Hashtag jest zapisywany jako
token \texttt{hashtag\_nazwa}, na przykład \texttt{\#MAGA} jako
\texttt{hashtag\_maga}. Takie rozwiązanie zachowuje dwie
informacje równocześnie: konkretną nazwę oraz fakt, że nazwa
została użyta jako wzmianka albo hashtag. Jest to korzystne dla
późniejszej interpretacji SHAP, ponieważ cecha \texttt{mentions\_foxnews}
jest bardziej jednoznaczna niż samo słowo \texttt{foxnews}. Dodatkowo
usunięto znacznik \texttt{rt}, który wskazuje, że wpis jest retweetem.

\section{Tokenizacja, stopwords i lematyzacja}

Po czyszczeniu tekst został podzielony na tokeny. Następnie
usunięto angielskie słowa, które występują bardzo często, ale zazwyczaj
wnoszą niewiele informacji o temacie, na przykład \emph{the},
\emph{and} lub \emph{of}. Nie zastosowano jednak bezwarunkowego
usunięcia wszystkich takich słów. Zachowano zaimki osobowe, przeczenia i
czasowniki modalne, między innymi \emph{i}, \emph{you}, \emph{we}, \emph{not},
\emph{never}, \emph{will} oraz \emph{should}. Są to słowa ważne dla stylu,
ponieważ mogą sygnalizować zwracanie się do odbiorców,
budowanie wspólnoty, deklarowanie pewności albo negowanie stanowiska przeciwnika.

Kolejnym krokiem była lematyzacja przy użyciu narzędzia WordNetLemmatizer.
Lemat to podstawowa forma słowa używana jako wspólny reprezentant jego
odmian (np. \emph{running}, \emph{ran} -> \emph{run}).
Zastosowanie lematyzacji zmniejsza liczbę odrębnych wariantów słów i
pozwala modelowi częściej rozpoznawać podobne znaczenia. Osobno usunięto
nazwę konta \texttt{realdonaldtrump}, ponieważ pojawia się ona często
jako samocytowanie i nie jest interesującym sygnałem językowym.

\begin{figure}[H]
\centering
\begin{subfigure}[t]{0.48\textwidth}
\centering
\autoplot{preprocessing_wordcloud.png}
\caption{Chmura słów po przetworzeniu tekstu.}
\label{fig:wordcloud}
\end{subfigure}
\hfill
\begin{subfigure}[t]{0.48\textwidth}
\centering
\autoplot{preprocessing_top_words.png}
\caption{Najczęściej występujące tokeny.}
\label{fig:top-words}
\end{subfigure}
\caption{Wizualne podsumowanie słownictwa po czyszczeniu i lematyzacji.}
\label{fig:preprocessing-viz}
\end{figure}

\chapter{Konstruowanie cech i reprezentacja tekstu}

\section{Cechy strukturalne i stylistyczne}

Słowa i frazy są głównym elementem projektu, ale tekst tweetów zawiera również cechy,
których nie opisuje sama lista słów. Z tego powodu obliczono 45 dodatkowych cech liczbowych.
Wykorzystano liczbę wykrzykników, znaków zapytania, słów zapisanych wielkimi literami,
słów, zdań i znaków. Dodano również liczbę hashtagów, wzmianek, adresów internetowych
i odnośników do mediów. Wielkie litery są na przykład wykorzystywane jako podkreślenie emocji.

Aby porównać krótkie i długie tweety, obliczono także wskaźniki względne.
Przykładowo \texttt{exclamation\_per\_word} oznacza liczbę wykrzykników
podzieloną przez liczbę słów powiększoną o jeden.
Cechy \texttt{ends\_with\_exclamation} oraz \texttt{ends\_with\_question} wskazują,
czy tweet kończy się odpowiednio wykrzyknikiem albo pytajnikiem.
Wskaźnik \texttt{max\_exclamation\_run} mierzy najdłuższy ciąg kolejnych
wykrzykników, a \texttt{shouting\_index} określa udział wielkich liter w całym wpisie.

Osobna grupa cech została zbudowana na podstawie list z określonym typem słownictwa. Obliczono
liczbę zaimków pierwszej osoby liczby pojedynczej i mnogiej, liczbę zwrotów w drugiej osobie,
czasowników modalnych i przeczeń. Obliczono również liczbę słów wzmacniających,
słów wyrażających pewność, superlatywów oraz wyraźnie negatywnych określeń.
Cechy te nie zastępują reprezentacji słów, ale umożliwiają modelowi uchwycenie
szerszych właściwości stylu, na przykład skłonności do używania kategorycznych
ocen albo bezpośrednich zwrotów do odbiorcy.

W projekcie uwzględniono również długość przeciętnego słowa, długość przeciętnego
zdania, wskaźnik czytelności Flesch Reading Ease, poziom Flesch--Kincaid Grade,
współczynnik powtórzeń słów oraz Type-Token Ratio. Type-Token Ratio jest
ilorazem liczby różnych słów i liczby wszystkich słów. Wysoka wartość
oznacza większą różnorodność słownictwa, a niska wartość wskazuje na powtarzalność.
Cechy czytelności należy interpretować ostrożnie, ponieważ zostały pierwotnie
opracowane dla dłuższych tekstów, podczas gdy tweet jest bardzo krótką formą wypowiedzi.

Jak stworzone cechy są powiązane z liczbą retweetów, pokazuje wykres \ref{fig:feature-correlation},
a zwłaszcza pierwsza kolumna którą podkreślono.
Widać, że niektóre cechy mają wyraźną korelację z liczbą retweetów.

\begin{table}[H]
\centering
\small
\caption{Grupy cech liczbowych wykorzystanych w końcowej macierzy danych.}
\label{tab:numerical-features}
\begin{tabularx}{\textwidth}{>{\raggedright\arraybackslash}p{0.26\textwidth}X}
\toprule
Grupa & Zastosowane kolumny \\
\midrule
Interpunkcja i długość & \texttt{exclamation\_count}, \texttt{question\_count}, \texttt{capslock\_words\_count}, \texttt{word\_count}, \texttt{sentence\_count}, \texttt{char\_count}. \\
Elementy platformy & \texttt{hashtag\_count}, \texttt{mention\_count}, \texttt{has\_url}, \texttt{has\_media\_link}, \texttt{url\_count}. \\
Intensywność zapisu & \texttt{exclamation\_per\_word}, \texttt{question\_per\_word}, \texttt{ends\_with\_exclamation}, \texttt{ends\_with\_question}, \texttt{max\_exclamation\_run}, \texttt{shouting\_index}. \\
Zaimek i retoryka & \texttt{first\_person\_singular\_count}, \texttt{first\_person\_plural\_count}, \texttt{second\_person\_count}, \texttt{modal\_count}, \texttt{negation\_count}, \texttt{emphatic\_word\_count}, \texttt{certainty\_word\_count}, \texttt{superlative\_word\_count}, \texttt{negative\_emphatic\_word\_count}. \\
Czytelność i słownictwo & \texttt{avg\_word\_length}, \texttt{avg\_sentence\_length}, \texttt{flesch\_reading\_ease}, \texttt{flesch\_kincaid\_grade}, \texttt{repetition\_ratio}, \texttt{lexical\_diversity\_ttr}. \\
Sentyment & \texttt{sent\_vader\_neg}, \texttt{sent\_vader\_neu}, \texttt{sent\_vader\_pos}, \texttt{sent\_vader\_compound}. \\
Liczby w tekście & \texttt{has\_number}, \texttt{number\_token\_count}, \texttt{has\_year\_mention}, \texttt{has\_large\_number}. \\
Tematy & \texttt{topic\_lda\_1}, \texttt{topic\_lda\_2}, \texttt{topic\_lda\_3}, \texttt{topic\_lda\_4}, \texttt{topic\_lda\_5}. \\
\bottomrule
\end{tabularx}
\end{table}

W czasie analizy stwierdzono, że \num{13118} tweetów, czyli około 30{,}3\%,
zawiera co najmniej jedną liczbę. Około 3{,}1\% wpisów (1339) zawiera rok w postaci czterech
cyfr, a około 0{,}2\% (93) zawiera inną liczbę co najmniej czterocyfrową.
Takie elementy mogą opisywać informacje dotyczące wyborów, dat, wyników
lub wartości ekonomicznych, dlatego zostały zachowane w zbiorze cech.

\begin{figure}[H]
\centering
\begin{subfigure}[t]{0.48\textwidth}
\centering
\autoplot{features_distributions.png}
\caption{Rozkłady cech stylistycznych.}
\label{fig:feature-distributions}
\end{subfigure}
\hfill
\begin{subfigure}[t]{0.48\textwidth}
\centering
\autoplot{features_numeric_tokens.png}
\caption{Związek liczb w tekście z logarytmem liczby retweetów.}
\label{fig:numeric-tokens}
\end{subfigure}
\caption{Wykresy z etapu tworzenia cech.}
\label{fig:feature-engineering-viz}
\end{figure}

\section{Analiza sentymentu VADER}

Dla każdego wpisu obliczono wynik sentymentu metodą VADER,
czyli \emph{Valence Aware Dictionary and sEntiment Reasoner}. Jest to słownikowo-regułowa
metoda przygotowana do tekstów z mediów społecznościowych \cite{hutto2014}. Zwraca ona cztery
podstawowe wartości: udział tonu negatywnego, neutralnego i pozytywnego oraz wynik złożony
\texttt{compound} w przybliżonym zakresie od $-1$ do $1$. Wartość ujemna sugeruje przewagę
tonu negatywnego, dodatnia przewagę tonu pozytywnego, a wartość bliska zeru wskazuje na
ton neutralny.

Do obliczania sentymentu wykorzystano oryginalną treść tweetów po usunięciu adresów i elementów technicznych. Zachowano natomiast interpunkcję, wielkie litery i emotikony, ponieważ VADER potrafi wykorzystywać te sygnały. Wartość \texttt{sentiment\_intensity}, równa wartości bezwzględnej wyniku \texttt{compound}, została utworzona do analizy pomocniczej. Do końcowej macierzy modelu weszły cztery podstawowe wyniki VADER: negatywny, neutralny, pozytywny i złożony.

\begin{figure}[H]
\centering
\autoplot{features_sentiment.png}
\caption{Rozkład kategorii sentymentu, rozkład wyniku złożonego VADER oraz logarytm liczby retweetów w zależności od sentymentu.}
\label{fig:sentiment}
\end{figure}

\section{Reprezentacja słów: CountVectorizer i TF-IDF}

Model klasyfikacyjny nie może bezpośrednio operować na ciągu znaków, dlatego tekst
zamieniono na wektor liczb. Pierwszą zastosowaną reprezentacją był CountVectorizer.
Dla każdego tweetu tworzy on wiersz macierzy, a dla każdego słowa lub frazy tworzy
kolumnę. Wartość w komórce oznacza liczbę wystąpień danego słowa albo frazy
w konkretnym tweecie. Taka reprezentacja jest nazywana Bag of Words, ponieważ
zachowuje informację o występowaniu słów, lecz nie zachowuje kolejności
w zdaniu ~\cite{CountVectorizer}.

Drugą reprezentacją był TF-IDF, czyli \emph{Term Frequency-Inverse Document Frequency}.
Pierwsza część tego skrótu, \emph{term frequency}, opisuje częstość danego terminu w
konkretnym dokumencie. Druga część, \emph{inverse document frequency}, obniża znaczenie
słów występujących w bardzo wielu dokumentach. Matematycznie jest to zapisane jako ~\cite{TFIDFequation}:
\[
\mathrm{TFIDF}(t,d) = \mathrm{TF}(t,d) \cdot \log\left(\frac{N}{\mathrm{DF}(t)+1}\right),
\]
gdzie $t$ oznacza termin, $d$ oznacza tweeta, $N$ oznacza liczbę tweetów,
a $\mathrm{DF}(t)$ oznacza liczbę tweetów zawierających termin $t$.
Dzięki temu często występujące słowa o małej wartości rozróżniającej otrzymują
niższą wagę, a słowa charakterystyczne dla mniejszej części wpisów otrzymują
wyższą wagę \cite{TFIDF}.

W obu reprezentacjach zastosowano jednoelementowe i dwuelementowe n-gramy.
Jednoelementowy n-gram to pojedyncze słowo, na przykład \emph{great}. Dwuelementowy
n-gram, nazywany bigramem, to para kolejnych słów, na przykład \emph{fake news}
albo \emph{thank you}. Uwzględnienie bigramów jest istotne, ponieważ znaczenie
frazy nie zawsze wynika z sumy znaczeń pojedynczych słów. Słownik ograniczono
do \num{10000} cech tekstowych. Dodatkowo odrzucono terminy występujące w
mniej niż pięciu tweetach oraz terminy obecne w ponad 95\% tweetów. Otrzymano
dwie macierze tekstowe o wymiarze \num{43352} na \num{10000}.

\section{Modelowanie tematów LDA}

W projekcie zastosowano również modelowanie tematów metodą LDA,
czyli \emph{Latent Dirichlet Allocation} \cite{blei2003}. Metoda ta
zakłada, że dokument może łączyć kilka ukrytych tematów, a każdy
temat jest opisywany przez rozkład słów. Zasadniczo algorythm LDA zakłada,
że dokumenty zostały wygenerowane poprzez losowe próbkowanie
tematów sprzed utworzenia dokumentu i próbuje dokonać inżynierii
wstecznej tego próbkowania, co przedstawione jest na stronie IBM \cite{IBM_LDA}.

Model LDA został dopasowany do surowych zliczeń słów z CountVectorizer,
ponieważ LDA operuje na reprezentacji bag of words.
Ustalono liczbę pięciu tematów. Dla każdego tweetu otrzymano
pięć liczb \texttt{topic\_lda\_1}-\texttt{topic\_lda\_5}, które
opisują stopień przypisania wpisu do poszczególnych tematów.
Wśród ważnych słów tematów pojawiły się między innymi
\emph{country}, \emph{china}, \emph{thank you}, \emph{democrat}, \emph{fake news}
oraz nazwy użytkowników \ref{fig:lda-topics}. Te cechy nie zastępują tekstowej reprezentacji TF-IDF,
ale dodają dodatkową informację o dominujących obszarach treści.

\begin{figure}[H]
\centering
\begin{subfigure}[t]{0.48\textwidth}
\centering
\autoplot{features_lda_topics.png}
\caption{Najważniejsze słowa dla pięciu tematów LDA.}
\label{fig:lda-topics}
\end{subfigure}
\hfill
\begin{subfigure}[t]{0.48\textwidth}
\centering
\autoplot{features_correlation_heatmap.png}
\caption{Macierz korelacji cech liczbowych i liczby retweetów.}
\label{fig:feature-correlation}
\end{subfigure}
\caption{Wizualizacja tematów oraz współzależności cech numerycznych.}
\label{fig:topics-correlation}
\end{figure}

\section{Łączenie i normalizacja cech}

Cechy tekstowe i cechy liczbowe mają różne skale. Przykładowo liczba znaków może
mieć wartość większą niż 100, wskaźnik sentymentu złożonego mieści się w pobliżu
przedziału od $-1$ do $1$, a cecha binarna ma wyłącznie wartości 0 i 1. Aby żadna
cecha nie dominowała wyłącznie z powodu swojej skali, wszystkie 45
cech liczbowych zostało przeskalowane metodą StandardScaler.
Po standaryzacji każda cecha ma średnią równą zero i odchylenie standardowe
równe jeden w zbiorze użytym do dopasowania skalera.

Następnie połączono kolumny tekstowe z kolumnami liczbowymi. Otrzymano dwie końcowe
rzadkie macierze:\\
\texttt{X\_tfidf\_features.npz} oraz \texttt{X\_count\_features.npz}.\\
Każda z nich ma \num{10045} kolumn, ponieważ zawiera \num{10000} cech tekstowych i
45 cech liczbowych. Macierz rzadka przechowuje przede wszystkim wartości
różne od zera. Jest to ważne, ponieważ pojedynczy tweet nie zawiera większości
słów obecnych w słowniku, a zapisywanie wszystkich zer byłoby bardzo nieefektywne
pamięciowo.

\chapter{Budowa i ocena modeli}

\section{Podział danych}

Dane zostały podzielone na zbiór treningowy i testowy w proporcji 80\% do 20\%.
Zbiór treningowy zawierał \num{34681} tweetów, a testowy \num{8671}.
Podział był stratyfikowany, co oznacza, że proporcje klas \emph{High}, \emph{Low} i \emph{Medium}
zostały zachowane w obu częściach danych.

Etykiety tekstowe zostały technicznie zakodowane jako liczby. Zastosowany LabelEncoder ~\cite{LabelEncoderGFG}
uporządkował je alfabetycznie, dlatego w aktualnym uruchomieniu kod \texttt{0} oznaczał
klasę \emph{High}, kod \texttt{1} klasę \emph{Low}, a kod \texttt{2} klasę \emph{Medium}.
To kodowanie nie oznacza, że klasa \emph{High} jest mniejsza od \emph{Low}.
Jest to wyłącznie techniczna reprezentacja kategorii.

\section{Porównywane modele}

Pierwszym punktem odniesienia był DummyClassifier ze strategią \emph{most frequent}.
Nie analizuje on żadnych cech i zawsze przewiduje klasę najczęściej występującą w
zbiorze treningowym. Jego wynik pokazuje minimalny poziom, który powinny wyraźnie
przekroczyć modele wykorzystujące tekst.~\cite{DummyClassifierGFG}

Następnie przetestowano regresję logistyczną.
Model oblicza dla każdej klasy wynik zależny od ważonej sumy cech, a następnie wybiera klasę o największym prawdopodobieństwie.
Jest to model liniowy, co oznacza, że wpływ pojedynczej cechy jest opisany przez współczynnik.
Ta właściwość ułatwia późniejszą interpretację SHAP.~\cite{LogisticRegression}

Kolejnym modelem był liniowy SVM (Support Vector Machine).
Model SVM poszukuje granic rozdzielających klasy w przestrzeni cech.~\cite{SVM}
W projekcie porównano regresję logistyczną i liniowy SVM zarówno dla
reprezentacji TF-IDF, jak i dla prostych zliczeń CountVectorizer.

Ostatnią grupą modeli były klasyfikatory XGBoost (eXtreme Gradient Boosting). XGBoost jest
metodą wzmacniania drzew decyzyjnych. Kolejne drzewa są budowane tak,
aby poprawiać błędy poprzednich drzew.~\cite{XGBoost} W projekcie
zastosowano 250 drzew o maksymalnej głębokości 4, szybkość uczenia
równą 0,05 oraz algorytm \texttt{hist}. XGBoost jest modelem nieliniowym,
dlatego stanowi porównanie dla modeli liniowych.

\section{Metryki jakości}

Jakość klasyfikatorów porównano za pomocą kilku metryk.
Accuracy, czyli dokładność, jest odsetkiem poprawnie sklasyfikowanych tweetów.
Balanced accuracy, czyli zrównoważona dokładność, jest średnią czułości obliczoną
osobno dla każdej klasy. Metryka ta jest szczególnie przydatna, gdy klasy mają
różną liczność.~\cite{BalancedAccuracy}

Dodatkowo wykorzystano macro F1 i weighted F1. Miara F1 łączy precyzję i czułość.
Precyzja odpowiada na pytanie, jaka część tweetów przypisanych przez model do danej
klasy rzeczywiście do niej należy. Czułość odpowiada na pytanie, jaką część
prawdziwych przykładów tej klasy model odnalazł. Macro F1 oblicza wynik F1
dla każdej klasy oddzielnie, a następnie traktuje wszystkie klasy równorzędnie.
Weighted F1 również liczy wynik dla każdej klasy, ale nadaje większą wagę
klasom liczniejszym. W tym projekcie jako główną metrykę wyboru modelu
przyjęto macro F1, ponieważ zależało nam na porównywalnej jakości
dla klas \emph{High}, \emph{Low} i \emph{Medium}.

\section{Wyniki benchmarku}

Tabela \ref{tab:benchmark} przedstawia wyniki uzyskane na zbiorze testowym.
Najlepszy wynik macro F1, równy 0,818, uzyskała regresja logistyczna
wykorzystująca TF-IDF oraz cechy liczbowe. Ten sam model osiągnął accuracy
równą 0,818 i balanced accuracy równą 0,818. Wyniki regresji logistycznej
i liniowego SVM dla TF-IDF były bardzo zbliżone, ale regresja logistyczna
była nieznacznie lepsza i została wybrana jako model końcowy.

\begin{table}[H]
\centering
\small
\caption{Wyniki modeli na zbiorze testowym. Najwyższa wartość macro F1 została wyróżniona pogrubieniem.}
\label{tab:benchmark}
\begin{tabularx}{\textwidth}{>{\raggedright\arraybackslash}Xcccc}
\toprule
Model & Accuracy & Balanced accuracy & Macro F1 & Czas uczenia [s] \\
\midrule
DummyClassifier, najczęstsza klasa & 0,335 & 0,333 & 0,167 & 0,0 \\
Regresja logistyczna, TF--IDF + cechy liczbowe & 0,818 & 0,818 & \textbf{0,818} & 45,2 \\
Liniowy SVM, TF--IDF + cechy liczbowe & 0,814 & 0,814 & 0,813 & 57,4 \\
Regresja logistyczna, Count + cechy liczbowe & 0,812 & 0,812 & 0,812 & 52,9 \\
Liniowy SVM, Count + cechy liczbowe & 0,786 & 0,786 & 0,785 & 29,8 \\
XGBoost, TF--IDF + cechy liczbowe & 0,794 & 0,794 & 0,793 & 79,4 \\
XGBoost, Count + cechy liczbowe & 0,793 & 0,793 & 0,792 & 8,2 \\
\bottomrule
\end{tabularx}
\end{table}

Różnica między najlepszym modelem a DummyClassifier jest bardzo duża.
Model bazowy uzyskał macro F1 równe 0,167, ponieważ praktycznie zawsze
wskazywał tę samą klasę. Reprezentacja TF-IDF okazała się lepsza od
prostych zliczeń słów zarówno dla regresji logistycznej, jak i dla liniowego SVM.
Oznacza to, że ważenie rzadziej występujących, bardziej charakterystycznych
terminów było tutaj użyteczne.

\begin{figure}[H]
\centering
\begin{subfigure}[t]{0.62\textwidth}
\centering
\autoplot{models_benchmark_metrics.png}
\caption{Porównanie metryk modeli.}
\label{fig:model-benchmark}
\end{subfigure}
\hfill
\begin{subfigure}[t]{0.34\textwidth}
\centering
\autoplot{models_training_time.png}
\caption{Czas uczenia modeli.}
\label{fig:model-time}
\end{subfigure}
\caption{Porównanie jakości i kosztu obliczeniowego modeli.}
\label{fig:model-benchmark-time}
\end{figure}

\section{Szczegółowa ocena modelu końcowego}

Dla wybranego modelu obliczono raport klasyfikacji. Klasa \emph{High} została
rozpoznana najlepiej, z wynikiem F1 równym 0,888. Klasa \emph{Low} uzyskała
F1 równe 0,826. Najtrudniejsza była klasa \emph{Medium}, dla której
F1 wyniosło 0,739. Jest to intuicyjne, ponieważ klasa środkowa leży
między dwiema pozostałymi klasami i może zawierać wpisy podobne zarówno
do tweetów o niskim, jak i wysokim zainteresowaniu.

\begin{table}[H]
\centering
\caption{Raport klasyfikacji dla regresji logistycznej z reprezentacją TF--IDF i cechami liczbowymi.}
\label{tab:classification-report}
\begin{tabular}{lrrrr}
\toprule
Klasa & Precyzja & Czułość & F1 & Liczba tweetów \\
\midrule
High & 0,899 & 0,878 & 0,888 & \num{2890} \\
Low & 0,815 & 0,838 & 0,826 & \num{2904} \\
Medium & 0,740 & 0,737 & 0,739 & \num{2877} \\
\midrule
Średnia macro & 0,818 & 0,818 & 0,818 & \num{8671} \\
\bottomrule
\end{tabular}
\end{table}

\begin{figure}[H]
\centering
\begin{subfigure}[t]{0.28\textwidth}
\centering
\autoplot{models_best_confusion_matrix.png}
\caption{Macierz pomyłek modelu końcowego.}
\label{fig:best-confusion}
\end{subfigure}
\hfill
\begin{subfigure}[t]{0.68\textwidth}
\centering
\autoplot{models_top3_confusion_matrices.png}
\caption{Macierze pomyłek trzech najlepszych modeli.}
\label{fig:top3-confusions}
\end{subfigure}
\caption{Ocena błędów klasyfikacji.}
\label{fig:confusions}
\end{figure}

\chapter{Analiza SHAP i dyskusja wyników}

\section{Cel analizy SHAP}

Wysoki wynik modelu nie odpowiada jeszcze na pytanie, dlaczego model podjął określoną decyzję.
Analiza SHAP (SHapley Additive exPlanations) służy do przypisania wkładu poszczególnych cech do wyniku
modelu \cite{SHAPGFG}.

Dla pojedynczego tweeta wynik modelu dla wybranej klasy można przybliżyć równaniem~\cite{SHAP}:
\[
f(\mathbf{x}) \approx \phi_0 + \sum_{j=1}^{p}\phi_j,
\]
gdzie $f(\mathbf{x})$ jest wynikiem modelu dla tweetu $\mathbf{x}$, $\phi_0$ jest wynikiem bazowym,
$p$ oznacza liczbę cech, a $\phi_j$ jest wkładem $j$-tej cechy.
Dodatnia wartość SHAP dla klasy \emph{High} oznacza, że dana cecha przesuwa wynik modelu w kierunku
tej klasy. Ujemna wartość oznacza, że przesuwa wynik w przeciwną stronę.~\cite{SHAP}

Ponieważ modelem końcowym była regresja logistyczna, zastosowano\\
\texttt{shap.LinearExplainer}.
Ten explainer jest dostosowany do modeli liniowych. Analizie poddano losową próbę \num{350}
tweetów ze zbioru testowego. Tło referencyjne zawierało \num{1000} przykładów treningowych,
jednak mechanizm maskujący SHAP użył podpróby 100 przykładów jako kompromisu pomiędzy
stabilnością wyniku a kosztem obliczeń. Dla każdej z \num{350} obserwacji obliczono
wkład \num{10045} cech dla trzech klas, co dało tablicę o wymiarze $(350, 10045, 3)$.

\section{Globalna ważność cech}

Globalna ważność cechy została obliczona jako średnia wartość bezwzględna SHAP,
uśredniona po analizowanych tweetach i klasach. Wysoka wartość oznacza,
że cecha często i silnie zmienia decyzje modelu, niezależnie od kierunku
tego wpływu. Dane zostały przedstawione w tabeli \ref{tab:shap-global} oraz wykresie \ref{fig:shap-global}.
Najważniejszą cechą globalną była liczba słów \texttt{word\_count}.
Następne pozycje zajęły liczba wzmianek \texttt{mention\_count},
wskaźnik czytelności Flesch Reading Ease, średnia długość zdania,
wynik Flesch-Kincaid Grade, liczba znaków, udział trzeciego tematu LDA,
obecność odnośnika do medium oraz liczba wykrzykników.

\begin{table}[H]
\centering
\caption{Najważniejsze cechy według średniej wartości bezwzględnej SHAP.}
\label{tab:shap-global}
\begin{tabular}{lr}
\toprule
Cecha & Średnia wartość $|\mathrm{SHAP}|$ \\
\midrule
\texttt{word\_count} & 0,531 \\
\texttt{mention\_count} & 0,494 \\
\texttt{flesch\_reading\_ease} & 0,457 \\
\texttt{avg\_sentence\_length} & 0,320 \\
\texttt{flesch\_kincaid\_grade} & 0,293 \\
\texttt{char\_count} & 0,275 \\
\texttt{topic\_lda\_3} & 0,261 \\
\texttt{has\_media\_link} & 0,246 \\
\texttt{exclamation\_count} & 0,235 \\
\texttt{max\_exclamation\_run} & 0,205 \\
\bottomrule
\end{tabular}
\end{table}

W całej analizie cechy ręcznie skonstruowane, nazywane dalej cechami inżynierskimi,
odpowiadały za około 65,6\% sumy globalnej ważności SHAP.
Tekstowe jedno- i dwuwyrazowe n-gramy odpowiadały za około 34,4\%.
\ref{fig:shap-feature-type} pokazuje ten udział.
Nie oznacza to, że reprezentacja słów była niepotrzebna. Przeciwnie,
porównanie modeli wykazało, że TF-IDF poprawia jakość klasyfikacji.
Wynik SHAP pokazuje natomiast, że po połączeniu obu rodzajów danych
model wykorzystywał także silne i łatwo mierzalne sygnały stylistyczne,
takie jak długość wpisu, liczba wzmianek oraz obecność mediów.

\begin{figure}[H]
\centering
\begin{subfigure}[t]{0.48\textwidth}
\centering
\autoplot{shap_global_importance.png}
\caption{Globalna ważność cech SHAP.}
\label{fig:shap-global}
\end{subfigure}
\hfill
\begin{subfigure}[t]{0.48\textwidth}
\centering
\autoplot{share_of_global_shap_importance_by_feature_type.png}
\caption{Udział cech tekstowych i inżynierskich w ważności SHAP.}
\label{fig:shap-feature-type}
\end{subfigure}
\caption{Globalna interpretacja modelu końcowego.}
\label{fig:shap-global-summary}
\end{figure}


\section{Ważność cech dla poszczególnych klas}

Dla klasy \emph{High} największe średnie wartości bezwzględne SHAP miały: liczba wzmianek,
czytelność Flesch Reading Ease, liczba słów, poziom Flesch-Kincaid Grade, liczba znaków,
trzeci temat LDA, liczba wykrzykników, obecność medium i długość ciągu wykrzykników.
\ref{fig:shap-high}

Dla klasy \emph{Low} szczególnie ważne były liczba słów, liczba wzmianek, czytelność,
długość zdania oraz obecność linku do medium.
\ref{fig:shap-low}

Dla klasy \emph{Medium} dominowały średnia długość zdania, obecność adresu internetowego,
liczba wykrzykników w stosunku do długości tweetu oraz długość tekstu.
\ref{fig:shap-medium}

Należy tutaj rozróżnić dwie informacje. Średnia wartość bezwzględna SHAP mówi,
że cecha jest ważna. Średnia wartość ze znakiem mówi, czy w analizowanej próbce
cecha przeciętnie przesuwała wynik w stronę danej klasy, czy od niej oddalała.
Jeżeli średnia wartość ze znakiem jest bliska zeru, nie oznacza to braku znaczenia.
Może oznaczać, że dla części tweetów cecha zwiększa wynik klasy, a dla innych go zmniejsza.
\ref{fig:shap-signed-by-class}

\begin{figure}[H]
\centering
\begin{subfigure}[t]{0.32\textwidth}
\centering
\autoplot{shap_class_high.png}
\caption{Klasa High.}
\label{fig:shap-high}
\end{subfigure}
\hfill
\begin{subfigure}[t]{0.32\textwidth}
\centering
\autoplot{shap_class_low.png}
\caption{Klasa Low.}
\label{fig:shap-low}
\end{subfigure}
\hfill
\begin{subfigure}[t]{0.32\textwidth}
\centering
\autoplot{shap_class_medium.png}
\caption{Klasa Medium.}
\label{fig:shap-medium}
\end{subfigure}
\caption{Najważniejsze cechy |SHAP| dla każdej klasy zainteresowania.}
\label{fig:shap-by-class}
\end{figure}


\begin{figure}[H]
\centering
\begin{subfigure}[t]{0.32\textwidth}
\centering
\autoplot{shap_signed_class_high.png}
\caption{Klasa High.}
\label{fig:shap-signed-high}
\end{subfigure}
\hfill
\begin{subfigure}[t]{0.32\textwidth}
\centering
\autoplot{shap_signed_class_low.png}
\caption{Klasa Low.}
\label{fig:shap-signed-low}
\end{subfigure}
\hfill
\begin{subfigure}[t]{0.32\textwidth}
\centering
\autoplot{shap_signed_class_medium.png}
\caption{Klasa Medium.}
\label{fig:shap-signed-medium}
\end{subfigure}
\caption{Najważniejsze cechy SHAP ze znakiem dla każdej klasy zainteresowania.}
\label{fig:shap-signed-by-class}
\end{figure}


\begin{figure}[H]
\centering
\begin{subfigure}[t]{0.35\textwidth}
\centering
\autoplot{shap_bar_class_high.png}
\caption{Klasa High.}
\label{fig:shap-bar-high}
\end{subfigure}
\hfill
\begin{subfigure}[t]{0.35\textwidth}
\centering
\autoplot{shap_beeswarm_class_high.png}
\caption{Klasa High.}
\label{fig:shap-beeswarm-high}
\end{subfigure}
\hfill
\begin{subfigure}[t]{0.35\textwidth}
\centering
\autoplot{shap_bar_class_low.png}
\caption{Klasa Low.}
\label{fig:shap-bar-low}
\end{subfigure}
\hfill
\begin{subfigure}[t]{0.35\textwidth}
\centering
\autoplot{shap_beeswarm_class_low.png}
\caption{Klasa Low.}
\label{fig:shap-beeswarm-low}
\end{subfigure}
\hfill
\begin{subfigure}[t]{0.35\textwidth}
\centering
\autoplot{shap_bar_class_medium.png}
\caption{Klasa Medium.}
\label{fig:shap-bar-medium}
\end{subfigure}
\hfill
\begin{subfigure}[t]{0.35\textwidth}
\centering
\autoplot{shap_beeswarm_class_medium.png}
\caption{Klasa Medium.}
\label{fig:shap-beeswarm-medium}
\end{subfigure}
\caption{Najważniejsze cechy |SHAP| dla każdej klasy zainteresowania z wykresami Beeswarm.}
\label{fig:shap-beeswarm-by-class}
\end{figure}

\section{Interpretacja cech tekstowych}

Wśród tekstowych n-gramów zauważalną ważność miały między \\
innymi \texttt{mentions\_realdonaldtrump}, \texttt{mentions\_foxnews} \\
oraz słowo \texttt{thanks}.

Szczególnie istotne jest zachowanie wzmianki w postaci \texttt{mentions\_nazwa}.
Gdyby usunąć wyłącznie znak \texttt{@}, model widziałby na przykład zwykły
token \texttt{cnn}. Dzięki prefiksowi \texttt{mentions\_} można odróżnić nazwę
używaną jako wzmianka od potencjalnego zwykłego słowa oraz łatwiej analizować
decyzje SHAP. Wynik ten wspiera decyzję o zachowaniu pełnych wzmianek
w tekście poddawanym wektoryzacji.

\section{Wyjaśnienia lokalne}

Poza analizą globalną przygotowano wyjaśnienia lokalne, czyli wyjaśnienia pojedynczych
tweetów. Dla każdego przykładu wykres \ref{fig:shap-local} pokazuje cechy najsilniej
przesuwające wynik modelu
w kierunku przewidzianej klasy oraz cechy działające w przeciwnym kierunku.


\begin{figure}[H]
\centering
\begin{subfigure}[t]{0.30\textwidth}
\centering
\autoplot{shap_local_high_example.png}
\caption{Lokalne wyjaśnienie przykładu przewidzianego jako High.}
\label{fig:local-high}
\end{subfigure}
\hfill
\begin{subfigure}[t]{0.30\textwidth}
\centering
\autoplot{shap_local_low_example.png}
\caption{Lokalne wyjaśnienie przykładu przewidzianego jako Low.}
\label{fig:local-low}
\end{subfigure}
\hfill
\begin{subfigure}[t]{0.30\textwidth}
\centering
\autoplot{shap_local_medium_example.png}
\caption{Lokalne wyjaśnienie przykładu przewidzianego jako Medium.}
\label{fig:local-medium}
\end{subfigure}
\caption{Przykładowe lokalne wyjaśnienia SHAP.}
\label{fig:shap-local}
\end{figure}

\chapter{Ograniczenia i odtwarzalność}

Najważniejszym ograniczeniem obecnej implementacji jest moment przygotowania cech.
CountVectorizer, TF-IDF, StandardScaler oraz model LDA zostały dopasowane do pełnego
zbioru danych przed wykonaniem podziału na zbiór treningowy i testowy.
Oznacza to, że słownik słów, wagi TF-IDF, parametry skalowania i tematy LDA
zostały obliczone z wykorzystaniem również tweetów, które później znalazły się w zbiorze testowym.
Czyli mamy przeciekiem danych, ponieważ część informacji o rozkładzie testowym trafia do etapu przygotowania danych.

Przeciek nie przekazuje modelowi bezpośrednio poprawnych etykiet klas, ale może nieznacznie
zawyżać raportowane wyniki. W projekcie przyjęto to uproszczenie ze względu na
przejrzystą organizację notebooków i zapis pełnych macierzy cech do plików.
W wersji stricte produkcyjnej należałoby najpierw podzielić surowe dane na zbiór
treningowy i testowy, następnie dopasować wszystkie przekształcenia wyłącznie
na danych treningowych, a później zastosować je do danych testowych.

Drugim ograniczeniem jest to, że liczba retweetów zależy od wielu czynników nieobecnych w
macierzy cech. Należą do nich czas publikacji, liczba obserwujących konto w danym okresie,
bieżące wydarzenia polityczne, ekspozycja medialna i działanie mechanizmów platformy.

Aby odtworzyć wyniki, należy pobrać dane ze strony Kaggle, umieścić pliki w katalogu
\texttt{data/raw}, utworzyć środowisko wirtualne i zainstalować wymagane biblioteki.
Następnie notebooki należy uruchomić w kolejności: \texttt{01\_EDA.ipynb},
\texttt{02\_preprocessing.ipynb}, \texttt{03\_features.ipynb}, \texttt{04\_models.ipynb}
oraz \texttt{05\_shap.ipynb}. Notebooki zapisują przetworzone dane, macierze cech,
model końcowy oraz wyniki analizy w odpowiednich katalogach projektu.

Całość kodu wraz z dokumentacją przygotowaną w latex'u oraz wszystkimi grafikami
jest dostępna na GitHubie pod adresem ~\cite{repository}.

\chapter{Wnioski}

W projekcie zbudowano system klasyfikujący tweety do trzech poziomów zainteresowania
mierzonego liczbą retweetów. Podstawą systemu była reprezentacja tekstu za pomocą
TF-IDF i CountVectorizer, rozszerzona o 45 cech opisujących formę i styl wypowiedzi.
Najlepszy rezultat uzyskała regresja logistyczna pracująca na reprezentacji TF-IDF
połączonej z cechami liczbowymi. Model osiągnął macro F1 równe 0,818 oraz accuracy
równe 0,818 na wydzielonym zbiorze testowym. Wynik ten był wyraźnie lepszy
od klasyfikatora bazowego, który zawsze wybierał najczęstszą klasę.

Porównanie modeli wskazało, że dla dużej, rzadkiej macierzy słów dobrze sprawdzają
się modele liniowe. TF-IDF okazał się lepszy niż proste zliczenia słów, co sugeruje,
że ważenie terminów charakterystycznych dla mniejszej części tweetów jest w
tym zadaniu przydatne. Analiza SHAP pokazała, że model korzysta nie tylko
z pojedynczych słów i bigramów, ale również z długości wpisu, liczby wzmianek,
cech czytelności, obecności linków do mediów, interpunkcji i tematów LDA.

Analiza SHAP umożliwiła przedstawienie tych zależności w sposób zrozumiały
i pozwoliła przejść od samego wyniku klasyfikacji do dyskusji o cechach,
które model rzeczywiście wykorzystuje.

\FloatBarrier
\printbibliography[title={Bibliografia}]

\end{document}
